% Chapter 7: inference results from the frozen R=5,000 archive.
\chapter{Inference Results}
\label{chap:inference-results}

Chapter~\ref{chap:inference-methodology} separated two potential sources of
discrepancy: the residual left by finite CG stopping and the approximation used
to calibrate the resulting moment statistic.  The results therefore begin with
the Krylov path itself.  Until $\mathcal T_{n,m}$ has approached its
direct-moment target $\mathcal S_n$, departures from the weighted chi-square
reference cannot be attributed primarily to calibration.  Once
this approximation has been assessed, power and confidence-set inversion show
what the test detects and which candidate slopes it retains.

The results concern the inference-specific CG--FPLS path rather than a new
comparison of all four estimators studied in Chapter~\ref{chap:estimation-results}.
This narrower focus is important: raw FPLS, Arnoldi FPLS, and FPCR enter the
thesis as estimation and prediction procedures, whereas the inferential theory
examined here is built around the late-stopped CG moment residual.

\section{Archive verification and reporting conventions}
\label{sec:inference-results-verification}

The frozen archive retains the configuration, replication-level statistics
and diagnostics, five numerical summary tables, and five vector figures.
Recomputing the summaries supported by the stored arrays recovers the
corresponding reported quantities to floating-point precision.  No stored
numerical array contains a missing or infinite value, and no replication is
removed from the analysis.

The four experimental branches retain the distinctions fixed in
Table~\ref{tab:inference-design}.  Null calibration evaluates
$\mathcal T_{n,70}$ at $n=100$ against a 50,000-draw continuous-design oracle
weighted chi-square reference.  The stopping experiment compares fixed indices,
the estimation-oriented adaptive variant, and the maximum-rank fit on a separate
Monte Carlo stream.  Power uses the direct statistic $\mathcal S_n$ and an
empirical null critical value for each model and sample size.  The inverted-set
illustration instead uses $\mathcal T_{n,10}$, oracle calibration, and an
analytic five-dimensional coefficient ellipsoid.  Each branch answers a
distinct inferential question and is therefore interpreted using its own
statistic and calibration.

Unless stated otherwise, probabilities are based on 5,000 replications per
model and design point.  Their reported Monte Carlo standard errors use
\eqref{eq:simulation-mcse}.  They condition on the realised simulated critical
value and therefore do not include uncertainty from estimating a 50,000-draw
oracle quantile or a 5,000-draw empirical null quantile.  This additional
uncertainty is not quantified in the archive.  The resulting error bars are
therefore conditional Monte Carlo summaries, not complete uncertainty
intervals for the calibration procedure.

\section{From estimation stopping to inferential stopping}
\label{sec:inference-results-stopping}

The first question is whether an estimation-oriented discrepancy threshold
also makes the CG fitting error negligible at the root-$n$ scale.  As specified
in Section~\ref{sec:inference-implementation}, the adaptive comparator retains
the earlier threshold formula and constants but caps its variance pilot at
20 components rather than 70 and uses a stabilised moment-residual path rather
than the short recurrence.  The results below concern this variant, rather
than an unchanged rerun of the Chapter~\ref{chap:estimation-results} procedure.
Figure~\ref{fig:inference-stopping-diagnostics} follows
the path through three linked quantities: the scaled residual
$\sqrt n\|e_m\|$, the difference
$|\mathcal T_{n,m}-\mathcal S_n|$, and the resulting null rejection frequency.

\begin{figure}[p]
    \centering
    \includegraphics[width=\textwidth]{stopping_diagnostics.pdf}
    \caption[Stopping diagnostics for inference]{Effect of the stopping index on CG--FPLS inference at $n=100$.
    The top row shows the median scaled moment residual $\sqrt n\|e_m\|$
    from \eqref{eq:inference-two-residuals}.  The dotted line is the median
    estimation-oriented stopping threshold for $\sqrt n\|e_m\|$, not the
    late-stopping target.
    The middle row shows the median absolute difference between the finite-$m$
    and direct statistics, and the bottom row shows rejection under the null
    using the oracle 5\% critical value.  Circles denote fixed $m$, stars the
    adaptive variant with a 20-component pilot cap, and diamonds the
    maximum-rank fit.  Each summary uses 5,000
    replications from the stopping-diagnostic stream.}
    \label{fig:inference-stopping-diagnostics}
\end{figure}

Table~\ref{tab:inference-stopping-selected} extracts the adaptive rule and two
fixed indices from the complete paths.  The adaptive rule selects a median of
one, two, and one components in Models~1--3.  More specifically, it selects
$m=1$ in 99.96\% of Model~1 samples, $m=2$ in 92.94\% of Model~2 samples, and
$m=1$ in 63.88\% of Model~3 samples.  These choices are consistent with the
rule's estimation objective, but the resulting null rejection frequencies are
52.70\%, 96.02\%, and 71.94\%.  These large discrepancies show that the
residual is not negligible in the first-order approximation and causes severe
finite-sample distortion in these designs; theoretically, it can enter the
statistic at first order rather than merely inducing a small
calibration error.

\begin{table}[htbp]
    \centering
    \scriptsize
    \caption[Selected stopping diagnostics]{Selected stopping diagnostics from the separate 5,000-replication
    experiment.  The residual and statistic-difference columns report medians;
    the final column reports the empirical rejection percentage.}
    \label{tab:inference-stopping-selected}
    \begin{tabular}{@{}llrrrr@{}}
        \toprule
        Model & Rule & Median $m$ & Median $\sqrt n\|e_m\|$
              & Median $|\mathcal T_{n,m}-\mathcal S_n|$ & Rejection (\%) \\
        \midrule
        Model 1 & Adaptive & 1  & 4.2452 & 7.6719 & 52.70 \\
                & $m=10$   & 10 & 0.2392 & 0.05770 & 6.16 \\
                & $m=70$   & 70 & 0.003113 & $1.05\times10^{-5}$ & 6.20 \\
        \addlinespace
        Model 2 & Adaptive & 2  & 9.4103 & 79.9799 & 96.02 \\
                & $m=10$   & 10 & 0.2711 & 0.07506 & 5.88 \\
                & $m=70$   & 70 & 0.003109 & $2.16\times10^{-5}$ & 6.00 \\
        \addlinespace
        Model 3 & Adaptive & 1  & 9.0848 & 69.7139 & 71.94 \\
                & $m=10$   & 10 & 0.3005 & 0.08947 & 5.20 \\
                & $m=70$   & 70 & 0.003113 & $1.01\times10^{-5}$ & 5.26 \\
        \bottomrule
    \end{tabular}
\end{table}

The failure is not caused by the adaptive rule missing its own threshold.  All
15,000 adaptive fits satisfy their realised estimation-oriented stopping
threshold.  On the $\sqrt n$ scale, the median thresholds are 11.82, 12.57,
and 15.25, all above the corresponding adaptive median residuals in
Table~\ref{tab:inference-stopping-selected}.  The rule therefore stops as
designed.  The problem is that an order-one threshold is compatible with
regularised estimation but not with the vanishing requirement in
\eqref{eq:inference-late-stopping-condition}.

Rejection frequency is not monotone in $m$.  For example, fixed $m=2$ yields
2.64\% rejection in Model~1 but 99.82\% in Model~2, while $m=3$ yields 3.12\%
in Model~3.  At $m=5$, the three rejection frequencies happen to lie between
4.40\% and 5.64\%, although the median statistic discrepancy still ranges from
0.458 to 1.203.  An apparently satisfactory size at one intermediate index can
therefore be accidental; the diagnostic evidence for late stopping is the
systematic decay of the residual and statistic difference, not monotonic
movement of the rejection rate toward 5\%.

By $m=10$, the median statistic discrepancies have fallen below 0.09 and the
three rejection frequencies are already close to their late-path values.  This
numerical proximity motivates the separate $m_C=10$ confidence-set
illustration in these designs, but it does not validate ten components as a
general inference rule.  At $m=70$, all three median scaled residuals are approximately 0.0031
and the median discrepancies are of order $10^{-5}$.  Continuing to maximum
rank reduces both quantities to numerical roundoff without changing any
rejection decision relative to $m=70$ in this stopping experiment.

The separately generated principal size stream provides a decision-level
check.  Across its 5,000 samples per model, the maximum absolute
differences between $\mathcal T_{n,70}$ and $\mathcal S_n$ are
$9.37\times10^{-5}$, $4.02\times10^{-4}$, and $1.18\times10^{-4}$, and the two
statistics produce identical decisions at every stored 1\%, 5\%, and 10\%
critical value.  Iteration error is therefore negligible for the recorded
decisions.  The remaining departure from nominal size combines finite-sample
calibration effects, the fixed-grid approximation described in
\eqref{eq:inference-grid-oracle-law}, and Monte Carlo error in the reference
distribution and critical value.  Increasing the CG index does not address
those sources of discrepancy.

\section{Finite-sample null law and empirical size}
\label{sec:inference-results-null}

Figure~\ref{fig:inference-null-calibration} compares the finite-sample law of
$\mathcal T_{n,70}$ with the continuous-design oracle weighted chi-square
reference from
\eqref{eq:inference-oracle-law}.  The distributional centres are close in all
three models, while the upper quantiles reveal the calibration discrepancy
most relevant to testing.

\begin{figure}[p]
    \centering
    \includegraphics[width=\textwidth]{null_calibration.pdf}
    \caption[Finite-sample null calibration]{Finite-sample null calibration of the late-stopped CG--FPLS test.
    Rows correspond to Models~1--3.  The left column overlays the empirical density of $\mathcal T_{n,70}$ from
    5,000 samples at $n=100$ with 50,000 draws from the continuous-design
    oracle weighted chi-square reference.  The right column plots finite-sample
    quantiles
    horizontally and reference quantiles vertically against the identity
    line.  Different density scales across models reflect the different null
    dispersions.  Agreement is strongest in the centre; points below the line
    in the upper tail indicate larger finite-sample quantiles.}
    \label{fig:inference-null-calibration}
\end{figure}

The empirical and reference means in
Table~\ref{tab:inference-null-summary} differ by only 0.05, 0.12, and 0.02 in
Models~1--3.  Dispersion is less accurately reproduced.  The finite-sample
variances exceed their simulated reference values by 11.6\%, 20.9\%, and
15.6\%, respectively.  Correspondingly, the empirical 95th percentiles are
3.47\%, 3.87\%, and 3.71\% above the oracle quantiles.  The Q--Q panels express
the same pattern: central points follow the diagonal closely, but the
finite-sample right tail is heavier than the continuous-design reference
predicts.  Because the sampled basis is not exactly orthonormal, this
comparison also includes a grid discrepancy that persists when $n$ increases
at fixed $T$.

\begin{table}[htbp]
    \centering
    \scriptsize
    \caption[Null distributions and empirical test size]{Null-distribution and empirical-size summary at $n=100$.
    Quantities labelled ``reference'' use the 50,000-draw continuous-design
    oracle weighted chi-square sample.  Rejection entries are percentages
    based on 5,000 finite-sample replications.}
    \label{tab:inference-null-summary}
    \begin{tabular}{@{}lrrrrrrr@{}}
        \toprule
        & \multicolumn{2}{c}{Mean}
        & \multicolumn{2}{c}{$q_{0.95}$}
        & \multicolumn{3}{c}{Rejection probability} \\
        \cmidrule(lr){2-3}\cmidrule(lr){4-5}\cmidrule(l){6-8}
        Model & Empirical & Reference & Empirical & Reference
              & 1\% level & 5\% level & 10\% level \\
        \midrule
        Model 1 & 8.590 & 8.543 & 15.582 & 15.059 & 1.20 & 5.90 & 11.18 \\
        Model 2 & 8.666 & 8.542 & 15.762 & 15.176 & 1.48 & 5.84 & 11.72 \\
        Model 3 & 14.273 & 14.253 & 27.333 & 26.355 & 1.46 & 6.10 & 10.72 \\
        \bottomrule
    \end{tabular}
\end{table}

Models~1 and 2 have the same population score variances and error variance, so
their direct-moment null laws and weighted chi-square limits are identical even
though their slopes differ.  Under the null, the slope cancels exactly from
the direct score average in
\eqref{eq:inference-null-score-average}.  Their small differences in the table
come from independent finite-sample and oracle simulation draws; the negligible
$m=70$ residual can add only a numerical difference to this comparison.
Model~3 has
a different covariance design and consequently a wider reference law, with a
95th percentile near 26.36 rather than 15.1.

At the principal 5\% level, the oracle-calibrated rejection frequencies are
5.90\%, 5.84\%, and 6.10\%, with conditional Monte Carlo standard errors of
approximately 0.33--0.34 percentage points.  The common direction of the
differences indicates mild liberal behaviour at $n=100$, although inference
about exact calibration would also need to account for uncertainty in the
simulated critical values.  The 1\% and 10\% results show the same broad
tendency, with magnitudes that vary across levels and models.
Figure~\ref{fig:app-inference-calibration-comparison} in
Appendix~\ref{app:inference-diagnostics} displays the 5\% result with
conditional Monte Carlo error bars.

The direction of the distortion agrees with the quantile comparison: the
oracle reference understates the finite-sample right tail and therefore rejects
slightly too often.  Because the $m=70$ and direct statistics are already
numerically indistinguishable, additional CG iterations would not correct this
feature.  For the power experiment, the analysis instead uses empirical null
critical values so that detectability is not confounded with this modest
oracle-calibration error.

\section{Power against fixed perturbations}
\label{sec:inference-results-power}

The power experiment tests the baseline candidate $b=\beta$ while generating
responses from $\beta_\delta(s)=\beta(s)+\delta s$.  Each non-zero value of
$\delta$ is a fixed alternative at the stated sample size; the curves should
not be read as a direct simulation of the root-$n$ local sequence in
\eqref{eq:inference-local-alternative}.  Moreover, the plotted statistic is
$\mathcal S_n$, not a finite-$m=70$ CG statistic.  Its critical value is the
empirical 95th percentile of a separate null distribution for each model and
$n$, as specified in Section~\ref{sec:inference-power}.

\begin{figure}[htbp]
    \centering
    \includegraphics[width=\textwidth]{power_curves.pdf}
    \caption[Empirical power under fixed alternatives]{Empirical power against
    $\beta_\delta(s)=\beta(s)+\delta s$.  Solid and dashed curves correspond to
    $n=100$ and $n=200$, respectively, and the dotted line marks 5\%.  Every
    point uses 5,000 alternative samples and the direct-moment statistic
    $\mathcal S_n$ with a model- and sample-size-specific empirical null
    critical value.  The alternatives use fixed $\delta$; the figure is not a
    finite-$m=70$ power calculation or a literal root-$n$ local-power
    experiment.}
    \label{fig:inference-power-curves}
\end{figure}

At $\delta=0$, rejection remains close to 5\% in all six model--sample-size
combinations, as shown in Table~\ref{tab:inference-power-selected}.  These
values are not forced to equal 5\%: the empirical critical value is obtained
from one null sample, whereas the plotted $\delta=0$ point uses another.  The
curves are approximately symmetric around zero, as expected for this two-sided
quadratic statistic under the symmetric Gaussian design.  Across the complete
grid, the largest difference between rejection at $\delta$ and $-\delta$ is
1.38 percentage points.  Because the two rejection rates come from separate
alternative samples, the standard error of their difference can be about one
percentage point; the observed asymmetry is modest relative to Monte Carlo
variation.

\begin{table}[htbp]
    \centering
    \scriptsize
    \caption[Selected empirical power values]{Selected empirical rejection probabilities from the power curves
    (percent).  The displayed non-zero values use the positive branch of the
    approximately symmetric alternative grid.}
    \label{tab:inference-power-selected}
    \begin{tabular}{@{}lrrrrrr@{}}
        \toprule
        & \multicolumn{3}{c}{$n=100$}
        & \multicolumn{3}{c}{$n=200$} \\
        \cmidrule(lr){2-4}\cmidrule(l){5-7}
        Model & $\delta=0$ & $\delta=0.25$ & $\delta=0.50$
              & $\delta=0$ & $\delta=0.25$ & $\delta=0.50$ \\
        \midrule
        Model 1 & 5.08 & 37.88 & 89.84 & 5.08 & 64.82 & 99.74 \\
        Model 2 & 5.24 & 36.70 & 89.34 & 4.58 & 64.80 & 99.54 \\
        Model 3 & 4.78 & 23.60 & 75.58 & 4.90 & 44.56 & 97.26 \\
        \bottomrule
    \end{tabular}
\end{table}

Doubling the sample size materially steepens every curve.  At
$\delta=0.25$, power rises from approximately 37\% to 65\% in Models~1 and 2
and from 23.6\% to 44.6\% in Model~3.  At $\delta=0.50$, Models~1 and 2 move
from about 90\% power to more than 99.5\%, while Model~3 rises from 75.6\% to
97.3\%.  These are fixed-alternative comparisons: increasing $n$ separates a
fixed covariance-weighted displacement more clearly from the null score
variation.

Models~1 and 2 again have nearly coincident curves.  They share the same
covariance operator and perturbation direction; because the null candidate is
the baseline $b=\beta$, subtracting $\langle X,b\rangle$ leaves
$\delta\langle X,s\rangle+\varepsilon$ in both models.  The small plotted
differences are Monte Carlo variation.  Model~3 is less
powerful over moderate values of $|\delta|$ for this particular direction.
That comparison reflects the joint effect of $K\Delta$ and the wider null law;
it is not evidence that Model~3 is uniformly less informative against every
functional alternative.

The power results are consistent with the theoretical consistency result and
demonstrate a clear finite-sample gain for the direct-moment target along
$\Delta(s)=s$.  They do not separately verify
the finite-$m$ CG approximation under every alternative, and empirical-null
calibration is not feasible from a single observed dataset.  Both limitations
matter when the simulation result is transferred to an application.

\section{Finite-dimensional inverted-set illustration}
\label{sec:inference-results-confidence}

The final experiment inverts $\mathcal T_{n,10}$ over the first five simulation
basis functions.  Figure~\ref{fig:inference-confidence-sets} displays the
resulting functional envelopes, while
Table~\ref{tab:inference-confidence-summary} separates three notions that a
plot alone cannot distinguish: non-rejection at the complete simulated slope,
acceptance of its five-term truncation, and the pointwise width of the
five-dimensional inverted set.

\begin{figure}[htbp]
    \centering
    \includegraphics[width=\textwidth]{confidence_sets.pdf}
    \caption[Five-dimensional inverted-set envelopes]{Finite-dimensional inverted-set illustration at $n=100$ using
    $m_C=10$, a five-dimensional analytic ellipsoid, and the oracle 5\%
    critical value.  The black curve is the complete true slope.  The grey
    region joins the pointwise medians of the lower and upper endpoints over
    5,000 separately realised envelopes; it is not one realised 95\%
    confidence band.  Dashed curves show the replication whose mean envelope
    width is closest to the median width.  All grey and dashed envelopes range
    only over the five-dimensional basis span.  Pointwise visual containment
    does not imply simultaneous membership in the coefficient ellipsoid.}
    \label{fig:inference-confidence-sets}
\end{figure}

\begin{table}[htbp]
    \centering
    \scriptsize
    \caption[Five-dimensional inverted-set summaries]{Summary of the analytic five-dimensional inverted sets.  The
    first two entries are acceptance frequencies in percent.  For each
    replication, width is averaged over the 200-point functional grid; the
    table then reports the mean and median of these replication-level widths.
    The final entry gives empty and unbounded frequencies in percent.  Every
    quantity uses 5,000 replications.}
    \label{tab:inference-confidence-summary}
    \begin{tabular}{@{}lccrrcc@{}}
        \toprule
        Model & \shortstack{Full-slope\\non-rejection}
              & \shortstack{Five-term\\acceptance}
              & \shortstack{Mean\\width}
              & \shortstack{Median\\width}
              & \shortstack{Median\\rank}
              & \shortstack{Empty/\\unbounded} \\
        \midrule
        Model 1 & 94.20 & 94.32 & 2.810 & 2.798 & 5 & 0/0 \\
        Model 2 & 94.20 & 94.12 & 2.824 & 2.817 & 5 & 0/0 \\
        Model 3 & 94.02 & 93.94 & 1.118 & 1.114 & 5 & 0/0 \\
        \bottomrule
    \end{tabular}
\end{table}

All 15,000 replication-specific inverted sets are non-empty, bounded, and have
rank five.  The
full-slope non-rejection frequencies are 94.20\%, 94.20\%, and 94.02\%, while
the five-term candidate is accepted in 94.32\%, 94.12\%, and 93.94\% of
samples.  The latter is the acceptance frequency of the five-term truncated
candidate in the candidate space underlying the display.  It should
nevertheless be interpreted with care: the data are generated from the
complete slope, whose non-zero tail is not treated as a nuisance component in
a separate projection model.

The full-slope values are paired with, but are not exactly one minus, the
principal size frequencies.  The confidence branch regenerates the same
samples and uses the same oracle critical values, but evaluates
$\mathcal T_{n,10}$ rather than $\mathcal T_{n,70}$.  The two stopping indices
change only 5, 6, and 6 of the 5,000 classifications in Models~1--3, producing
94.20\%, 94.20\%, and 94.02\% non-rejection at $m_C=10$, compared with
94.10\%, 94.16\%, and 93.90\% obtained by complementing the reported
$m=70$ sizes.  The closeness is consistent with the stopping diagnostics, but
the shared samples make this a paired check rather than independent coverage
evidence.

Models~1 and 2 have almost identical mean widths, about 2.81 and 2.82,
consistent with their shared predictor covariance geometry.  Model~3's mean
width is 1.12, approximately 60\% smaller in this five-dimensional
representation.
The stronger leading score variances provide more information for the
displayed coefficient directions, even after allowing for Model~3's larger
critical value.  This is a statement about the chosen leading subspace and
parameterisation, not uniformly narrower uncertainty for the entire
infinite-dimensional slope.

The display is descriptive: neither the pointwise median envelope nor the
representative boundaries constitute a realised simultaneous confidence band.
Moreover, the interpretation of the $m_C=10$ inversion remains specific to
these designs and conditional on the finite-iteration residual being
sufficiently small.

Taken together, the experiments identify a coherent sequence.  The
estimation-oriented adaptive variant stops with a moment residual that is
large on the root-$n$ scale and, in these designs, severely distorts the test.
Continuing the stabilised Krylov path makes that algorithmic error negligible;
at $m=70$, the test calibrated against the continuous-design oracle reference
is only mildly liberal at $n=100$.  The remaining discrepancy includes
finite-sample, grid, and reference-simulation effects.  Empirically calibrated
direct-moment power then increases sharply with sample size, and inversion yields non-empty,
bounded, full-rank finite-dimensional sets in all three controlled designs.

The scope is correspondingly limited.  The evidence uses dense Gaussian
curves, homoskedastic errors, a single perturbation direction, and either
oracle or simulation-specific empirical calibration.  No stored result
establishes the performance of a feasible plug-in or bootstrap calibration on
observed data.  Chapter~\ref{chap:spectroscopy-application} therefore shifts
from oracle-assisted inferential benchmarking to prediction with the four
regularisation procedures on spectroscopy curves.  A formal test in that
application would require a feasible calibration whose finite-sample behaviour
has not been evaluated here.  The principal empirical comparison is kept
concise in Chapter~\ref{chap:spectroscopy-application}; the complete
preprocessing, tuning, stability, and sensitivity record is reserved for
Appendix~\ref{app:spectroscopy-supplement}.
